\documentclass[11pt,letterpaper]{article}
\usepackage[letterpaper,margin=0.5in]{geometry}
\usepackage[utf8]{inputenc}
\usepackage{mdwlist}
\usepackage[default]{lato}
\usepackage[T1]{fontenc}
\usepackage{textcomp}
\usepackage{fontawesome}
\usepackage{enumitem}
\DeclareFontFamily{U}{fontawesomeOne}{}
\DeclareFontShape{U}{fontawesomeOne}{m}{n}
{<-> FontAwesome--fontawesomeone}{}
\DeclareRobustCommand\FAone{\fontencoding{U}\fontfamily{fontawesomeOne}\selectfont}
\pagestyle{empty}
\setlength{\tabcolsep}{0em}

% indentsection style, used for sections that aren't already in lists
% that need indentation to the level of all text in the document
\newenvironment{indentsection}[1]%
{\begin{list}{}%
{\setlength{\leftmargin}{#1}}%
     \item[]
}
{\end{list}}

% opposite of above; bump a section back toward the left margin
\newenvironment{unindentsection}[1]%
{\begin{list}{}%
{\setlength{\leftmargin}{-0.5#1}}%
\item[]%
}
{\end{list}}

% format two pieces of text, one left aligned and one right aligned
\newcommand{\headerrow}[2]
{\begin{tabular*}{\linewidth}{l@{\extracolsep{\fill}}r}
#1 &
#2 \\
\end{tabular*}}

% make "C++" look pretty when used in text by touching up the plus signs
\newcommand{\CPP}
{C\nolinebreak[4]\hspace{-.05em}\raisebox{.22ex}{\footnotesize\bf ++}}

% and the actual content starts here
\begin{document}

\begin{center}
	{\LARGE \textbf{Daniel Pizarro}}\\
	Santiago, Chile
	\vspace{0.05cm}
	\\
	\raisebox{-0.2\height} {\Large \faPhoneSquare} \ \  (000) 000-0000 \hfill\raisebox{-0.2\height}{\Large \faEnvelopeSquare} \ \ example@example.com \hfill \raisebox{-0.2\height}{\Large \faGithubSquare} \ \ github.com/example-profile \hfill \raisebox{-0.2\height}{\Large \faLinkedinSquare} \ \ linkedin.com/in/example
\end{center}

\hrule
\vspace{-1em}
\subsection*{\Large Experiencia}

\renewcommand\labelitemi{}
\renewcommand\labelitemii{$\bullet$}
\begin{itemize}[leftmargin=1em]
	\parskip=0.1em
		
		
	\item
	      \headerrow
	      {\textbf{SKY Airline}}
	      {\textbf{Santiago, Chile}}
	      \headerrow
	      {\emph{Desarrollador}}
	      {\emph{Diciembre 2021 -- Present}}
	      \begin{itemize*}
                                           \item - Formé parte del proyecto estratégico de migración del host de reservas de Radixx a SabreSonic, como parte de la transformación tecnológica de SKY y del desarrollo de sus nuevos canales digitales con miras a 2030. Lideré investigaciones técnicas sobre las interfaces de Sabre, incluyendo REST y Web Services, trabajando directamente con sus equipos para comprender los procesos, analizar problemas de integración y resolver dudas técnicas. También apoyé a los equipos de negocio en el entendimiento de las interfaces y los procesos involucrados, y participé en el cutover hacia la nueva plataforma.
                                      \item Lideré una iniciativa de desarrollo de un microservicio para centralizar la configuración de horarios operacionales asociados al proceso de check-in, incluyendo la hora de presentación de los pasajeros, la apertura y el cierre de los counters. La solución contempló reglas según rutas nacionales e internacionales, persistencia en base de datos y una interfaz de administración para los equipos operacionales de SKY.
                                        \item Desarrollé componentes de interfaz para un sistema de diseño documentado en Storybook, trabajando en conjunto con los equipos de UI/UX y Marketing para alinear su implementación con los diseños y lineamientos de marca.
                                \item Participé en el desarrollo de un servicio de check-in automático que procesaba los vuelos y ejecutaba el check-in según reglas de negocio y configuraciones por país, aeropuerto, tipo de pasajero y reserva. También participé en la coordinación con los aeropuertos para habilitar el servicio y en su marcha blanca en Osorno.
                      \item Participé en el desarrollo de un servicio de check-in automático que procesaba los vuelos y ejecutaba el check-in según reglas de negocio y configuraciones por país, aeropuerto, tipo de pasajero y reserva. También participé en la coordinación con los aeropuertos para habilitar el servicio y en su marcha blanca en Osorno.
            \item Participé en la migración a microservicios del flujo de generación de tarjetas de embarque y de los procesos encargados de enviarlas a los pasajeros a través de distintos canales, incluyendo correo electrónico y WhatsApp.
    	      \item Participé en la migración de funcionalidades del PSS Gateway, un monolito integrado con el host de reservas Radixx, hacia una arquitectura de microservicios. Trabajé en el caso de uso encargado de actualizar y persistir información de contacto y documentación de los pasajeros requerida para el check-in y embarque, según las regulaciones del país de destino.
	      \end{itemize*}
	      	      
	\item
	      \headerrow
	      {\textbf{Tubesoft}}
	      {\textbf{Santiago, Chile}}
	      \headerrow
	      {\emph{Desarrollador}}
	      {\emph{Marzo 2021 -- Noviembre 2021}}
	      \begin{itemize*}
	      	\item Participé como desarrollador para el cliente Liquidos.cl, enfocado en la mantención y evolución de su canal de ventas digital, con especial énfasis en los componentes de carrito de compras, checkout y procesamiento de pagos.
    	      \item Brindé soporte y realicé mejoras evolutivas sobre los flujos de pago del e-commerce, incluyendo la integración con Transbank (Webpay) y los componentes asociados al carrito de compras y checkout.
Además refactoricé y realicé mejoras sobre desarrollos existentes para facilitar su lectura y mantenimiento, reorganizando el código y separando responsabilidades.
    	      \item Brindé soporte y realicé mejoras evolutivas sobre los flujos de pago del e-commerce, incluyendo la integración con Transbank (Webpay) y los componentes asociados al carrito de compras y checkout.
Además refactoricé y realicé mejoras sobre desarrollos existentes para facilitar su lectura y mantenimiento, reorganizando el código y separando responsabilidades.
   	      \item Realicé ajustes en el logging para mejorar el seguimiento de la ejecución y facilitar la identificación de errores.
Lideré una iniciativa técnica para incorporar pruebas unitarias en el frontend, promoviendo la calidad y confiabilidad del código.
	      \end{itemize*}
	      	      
\end{itemize}

\hrule
\vspace{-1em}
\subsection*{\Large Habilidades}

\hyphenpenalty=1000
\begin{itemize}[leftmargin=1em,noitemsep]
	\item \textbf{Lenguajes:}
	      Javascript, HTML/CSS, Java, \LaTeX, Python, SQL
	\item \textbf{Herramientas y frameworks:}
	      Claude Code, Opencode, AWS (EC2, Lambda, RDS, S3), Microsoft (Azure Blob Storage, Azure Data Factory), NestJS, React.js, React Native.
\end{itemize}

\hrule
\vspace{-1em}
\subsection*{\Large Educación}

\begin{itemize}[leftmargin=1em]
	\parskip=0.1em
		
	\item
	      \headerrow
	      {\textbf{Instituto Profesional Duoc UC}}
	      {\textbf{Santiago, Chile}}
	      \headerrow
	      {\emph{Analista programador} }
	      {\emph{Marzo 2019 -- Julio 2021}}
	      	      
\end{itemize}

\end{document}
